% GENERATED FILE. Edit canonical lessons, then run npm run generate:books.
% canonical-source-hash: fnv1a64:fbf18e9ea20bfaab
% canonical-lessons: PT-C335-opiniao, PT-C335-convencido, PT-C335-ponto-de-vista, PT-C335-opinar, PT-C335-sustentar

\chapter{Na minha opinião}
\label{ch:pt-na-minha-opiniao}
\hlchaptermodality{335}

\begin{chapteropening}
\textbf{By the end of this chapter:} \emph{I can give my opinion and my point of view, say that I am convinced and maintain an idea.}

Complete the last lesson of chapter 335: I can give my opinion and my point of view, say that I am convinced and maintain an idea.
\end{chapteropening}

\section[a opinião]{a opinião — the opinion}

\label{lesson:PT-C335-opiniao}



\begin{quote}
Before the new one: say the Portuguese for a jungle, then the Portuguese for a slipper.
\end{quote}

\subsection*{You'll want to know: a opinião}
\textbf{a opinião} — \textquotedblleft{}the opinion\textquotedblright{}.

\begin{quote}
Cada pessoa tem a sua opinião.
\end{quote}

\begin{quote}
\emph{Every person has their own opinion.}
\end{quote}

\emph{Na minha opinião} opens a view; \emph{mudar de opinião} is to change your mind.

\begin{cousinweb}
\emph{opinião} comes from Latin \emph{opīniō, opīniōnis} \textquotedblleft{}opinion, supposition, belief\textquotedblright{} (Lewis \& Short), through Old Portuguese \emph{opynion} (EtymDB); it is built on the verb \emph{opīnārī} \textquotedblleft{}to suppose\textquotedblright{}.
\end{cousinweb}

\begin{grammarlens}[title={What you've built}]
One more word for telling, explaining and describing at B1.
\end{grammarlens}

\subsection*{Your turn}
\begin{itemize}
\raggedright
  \item \emph{Say it:} \emph{a opinião}
  \item \emph{Say it:} \emph{a opinião}, once more
  \item \emph{Say it:} a sentence of your own with \emph{a opinião}
  \item \emph{Recall:} say the Portuguese for a jungle, then the Portuguese for a slipper, then say \emph{a opinião} again
\end{itemize}

\begin{tcolorbox}[breakable,colback=teal!4,colframe=teal!35!black,title={Before you move on}]
What does a opinião mean? (\textquotedblleft{}The opinion\textquotedblright{}.) Say it once more.
\end{tcolorbox}
\section[convencido]{convencido — convinced, sure}

\label{lesson:PT-C335-convencido}



\begin{quote}
Before the new one: say the Portuguese for a slipper, then the Portuguese for the opinion.
\end{quote}

\subsection*{You'll want to know: convencido}
\textbf{convencido} — \textquotedblleft{}convinced, sure\textquotedblright{}.

\begin{quote}
Estou convencido de que é a melhor solução.
\end{quote}

\begin{quote}
\emph{I am convinced that it is the best solution.}
\end{quote}

\emph{Estar convencido de que…} — to be convinced that… A woman says \emph{convencida}. It is the past participle of \emph{convencer}.

\begin{cousinweb}
\emph{convencido} is the past participle of \emph{convencer}, which comes from Latin \emph{convincere} \textquotedblleft{}to overcome, to convict of error\textquotedblright{} (Lewis \& Short), built on \emph{vincere} \textquotedblleft{}to conquer\textquotedblright{}.
\end{cousinweb}

\begin{grammarlens}[title={What you've built}]
One more word for telling, explaining and describing at B1.
\end{grammarlens}

\subsection*{Your turn}
\begin{itemize}
\raggedright
  \item \emph{Say it:} \emph{convencido}
  \item \emph{Say it:} \emph{convencido}, once more
  \item \emph{Say it:} a sentence of your own with \emph{convencido}
  \item \emph{Recall:} say the Portuguese for a slipper, then the Portuguese for the opinion, then say \emph{convencido} again
\end{itemize}

\begin{tcolorbox}[breakable,colback=teal!4,colframe=teal!35!black,title={Before you move on}]
What does convencido mean? (\textquotedblleft{}Convinced, sure\textquotedblright{}.) Say it once more.
\end{tcolorbox}
\section[o ponto de vista]{o ponto de vista — the point of view}

\label{lesson:PT-C335-ponto-de-vista}



\begin{quote}
Before the new one: say the Portuguese for the opinion, then the Portuguese for convinced.
\end{quote}

\subsection*{You'll want to know: o ponto de vista}
\textbf{o ponto de vista} — \textquotedblleft{}the point of view\textquotedblright{}.

\begin{quote}
Do meu ponto de vista, a ideia é boa.
\end{quote}

\begin{quote}
\emph{From my point of view, the idea is good.}
\end{quote}

The plural is \emph{os pontos de vista}. \emph{Do ponto de vista de…} — from the point of view of…

\begin{cousinweb}
\emph{ponto} is inherited from Latin \emph{punctum} \textquotedblleft{}a prick, a point\textquotedblright{} (Lewis \& Short), and \emph{vista} is built on \emph{ver} \textquotedblleft{}to see\textquotedblright{}, inherited from Latin \emph{vidēre}: the spot you see things from.
\end{cousinweb}

\begin{grammarlens}[title={What you've built}]
One more word for telling, explaining and describing at B1.
\end{grammarlens}

\subsection*{Your turn}
\begin{itemize}
\raggedright
  \item \emph{Say it:} \emph{o ponto de vista}
  \item \emph{Say it:} \emph{o ponto de vista}, once more
  \item \emph{Say it:} a sentence of your own with \emph{o ponto de vista}
  \item \emph{Recall:} say the Portuguese for the opinion, then the Portuguese for convinced, then say \emph{o ponto de vista} again
\end{itemize}

\begin{tcolorbox}[breakable,colback=teal!4,colframe=teal!35!black,title={Before you move on}]
What does o ponto de vista mean? (\textquotedblleft{}The point of view\textquotedblright{}.) Say it once more.
\end{tcolorbox}
\section[opinar]{opinar — to give an opinion, to express a view}

\label{lesson:PT-C335-opinar}



\begin{quote}
Before the new one: say the Portuguese for convinced, then the Portuguese for the point of view.
\end{quote}

\subsection*{You'll want to know: opinar}
\textbf{opinar} — \textquotedblleft{}to give an opinion, to express a view\textquotedblright{}.

\begin{quote}
Toda a gente quis opinar sobre o novo jardim.
\end{quote}

\begin{quote}
\emph{Everyone wanted to give an opinion on the new garden.}
\end{quote}

\emph{Opinar sobre} — to give a view on. A little formal; in speech people say \emph{dar a opinião}.

\begin{cousinweb}
\emph{opinar} is a learned borrowing from Latin \emph{opīnārī} \textquotedblleft{}to be of opinion, to suppose, to believe\textquotedblright{} (EtymDB, Lewis \& Short).
\end{cousinweb}

\begin{grammarlens}[title={What you've built}]
One more word for telling, explaining and describing at B1.
\end{grammarlens}

\subsection*{Your turn}
\begin{itemize}
\raggedright
  \item \emph{Say it:} \emph{opinar}
  \item \emph{Say it:} \emph{opinar}, once more
  \item \emph{Say it:} a sentence of your own with \emph{opinar}
  \item \emph{Recall:} say the Portuguese for convinced, then the Portuguese for the point of view, then say \emph{opinar} again
\end{itemize}

\begin{tcolorbox}[breakable,colback=teal!4,colframe=teal!35!black,title={Before you move on}]
What does opinar mean? (\textquotedblleft{}To give an opinion, to express a view\textquotedblright{}.) Say it once more.
\end{tcolorbox}
\section[sustentar]{sustentar — to maintain, to claim; to support (a family, a weight)}

\label{lesson:PT-C335-sustentar}



\begin{quote}
Before the new one: say the Portuguese for the point of view, then the Portuguese for to give an opinion.
\end{quote}

\subsection*{You'll want to know: sustentar}
\textbf{sustentar} — \textquotedblleft{}to maintain, to claim; to support (a family, a weight)\textquotedblright{}.

\begin{quote}
Ele sustenta que o comboio é mais rápido do que o carro.
\end{quote}

\begin{quote}
\emph{He maintains that the train is faster than the car.}
\end{quote}

Also to support a family: \emph{sustentar a família}.

\begin{cousinweb}
\emph{sustentar} is a learned borrowing from Latin \emph{sustentāre} \textquotedblleft{}to hold up, to support, to sustain\textquotedblright{} (EtymDB, Lewis \& Short), a frequentative of \emph{sustinēre}, built from \emph{sub} \textquotedblleft{}under\textquotedblright{} and \emph{tenēre} \textquotedblleft{}to hold\textquotedblright{}.
\end{cousinweb}

\begin{grammarlens}[title={What you've built}]
One more word for telling, explaining and describing at B1.
\end{grammarlens}

\subsection*{Your turn}
\begin{itemize}
\raggedright
  \item \emph{Say it:} \emph{sustentar}
  \item \emph{Say it:} \emph{sustentar}, once more
  \item \emph{Say it:} a sentence of your own with \emph{sustentar}
  \item \emph{Recall:} say the Portuguese for the point of view, then the Portuguese for to give an opinion, then say \emph{sustentar} again
\end{itemize}

\begin{tcolorbox}[breakable,colback=teal!4,colframe=teal!35!black,title={Before you move on}]
What does sustentar mean? (\textquotedblleft{}To maintain, to claim; to support (a family, a weight)\textquotedblright{}.) Say it once more.
\end{tcolorbox}
